\label{ch:info_als_energie}
% Quelldokument: 255

\epigraph{Information ist topologisch eingefrorene kinetische Energie.}{Dok.\,255}

\section{Zwei Energietypen in \texorpdfstring{$\tilde{T} \cdot m = 1$}{T·m = 1}}

Die Grundrelation der FFGFT $\tilde{T} \cdot m = 1$ erzwingt eine natürliche
Zweiteilung der Energie: statische Energie als topologisch eingeschlossene
Kinematik und freie kinetische Energie als Bewegung durch $T^4$.

\textbf{Statische Energie:} An Ruhemasse gebundene Energie $E_0 = mc^2$
entspricht topologisch stabilen Windungszahlen auf $T^4$. Sie ist diskret, trägt
nicht zur thermodynamischen Entropie bei, und bleibt erhalten, solange die
Topologie intakt ist.

\textbf{Dynamische kinetische Energie:} Freie Bewegungsenergie --- bei
masselosen Teilchen die einzige Energieform ($E = hf$), bei massebehafteten
Teilchen der Anteil über die Ruhemasse hinaus. Sie ist kontinuierlich,
thermodynamisch und entropieproduktiv.

\section{Information als eingefrorene Kinematik}

Aus dieser Zweiteilung folgt ohne zusätzliche Postulate:

\begin{keyresult}
Information ist topologisch eingefrorene kinetische Energie. Eine
Windungskonfiguration $(n_\theta, n_\varphi)$ auf $T^4$ ist eine
kinetische Energie, die durch die Topologie des Torus gegen Dissipation
gesichert ist. Sie kann nicht verloren gehen, ohne dass die Topologie
bricht --- und topologische Übergänge sind in FFGFT quantisiert.
\end{keyresult}

Konsequenz: Landauers Prinzip beschreibt den Preis für das
\emph{Auftauen} eingefrorener Kinematik. Das Löschen eines Bits ist die
erzwungene Überführung einer topologisch gesicherten Konfiguration in eine
andere --- der dabei freigesetzte Phasenrauminhalt geht ans Bad.

\section{Quantisierung der Information}

Die Diskretheit der Windungszahlen erzwingt die Quantisierung der
Information. Es gibt kein halbiertes Bit --- die Windungszahl ist ganzzahlig.
Das ist kein Postulat, sondern eine direkte Konsequenz der Torus-Topologie.

Das Windungsquant $\Delta w = 1$ ist damit die fundamentale
Informationseinheit der FFGFT --- kleiner geht nicht, weil die Topologie keine
nicht-ganzzahligen Windungszahlen kennt.

\section{M.\,M.~Vopson aus FFGFT-Perspektive}

Vopsons Äquivalenz \textbf{[V1]} $E_{\text{info}} = mc^2_{\text{Bit}}$ lässt sich
aus FFGFT herleiten --- aber nur für den Spezialfall einer bestimmten Skala.
Die Energie eines Windungsquants auf Skala $L$ ist $\hbar c / L$, und die
zugehörige äquivalente Masse ist $m = \hbar c / (Lc^2) = \hbar / (Lc)$.
Das ist keine universelle Bitmasse --- sie ist skalensensitiv.

Vopsons \textbf{[V1]} Ergebnis ist der Grenzfall $L \to L_P$ (Planck-Skala):
dann gilt $m_{\text{bit}} \to m_P$, die Planck-Masse. Matzkes Ableitung
$m_{\text{bit}} = m_P \ln 2 / (2\pi)$ ist der Landauer-Spezialfall bei
Planck-Temperatur. Beide sind Grenzfälle der FFGFT-Formel, keine
universellen Gesetze.

\section{Der Casimir-Effekt als Beleg}

Der Casimir-Effekt zeigt, dass das Vakuum selbst Information trägt:
Die Vakuum-Energiedichte bei der charakteristischen $\xi$-Länge $L_\xi$
stimmt mit der Casimir-Energiedichte überein. Das ist kein Zufall ---
es ist die FFGFT-Vorhersage, dass $\xi$ sowohl die geometrische Skala
der Torus-Windungen als auch die Vakuum-Energieskala festlegt (Dok.\,025).
